\section{下降引理在非凸地形中只保证接近驻点}
\label{sec:09-Convex-Optimization/08-Nonconvex-and-Stochastic-Optimization/01}

把一段跑熟了的凸优化训练代码搬到神经网络上，改动往往少得可疑：损失函数换掉，模型换掉，更新那一行 \(x\leftarrow x-\eta\nabla f(x)\) 一个字符没动。跑起来损失也确实在降，梯度范数也确实在掉。于是很自然地沿用旧的收尾话术——``梯度范数到 \(10^{-5}\)，收敛了''。

这句话在凸的时候是一个结论，在非凸的时候是一句描述。中间丢掉的东西值得逐条清点，因为清点完你会发现：算法的行为没有退化，退化的是你能从它的行为里读出的信息。

\subsection{下降不等式从头到尾没有用到凸性}

\hyperref[sec:09-Convex-Optimization/06-Optimization-Algorithms/01]{``梯度下降真正难在步长''}里推过下降引理。回看那段推导的每一步：Taylor 展开取到一阶，余项用梯度的 Lipschitz 常数 \(L\) 控制，代入更新式，配方。四步里没有一步引用过凸性。

于是这个结论在非凸地形上原样成立。设 \(f\) 是 \(L\)-光滑的，更新为 \(x_{k+1}=x_k-\eta\nabla f(x_k)\)，则

\[
f(x_{k+1})\le f(x_k)-\eta\left(1-\frac{L\eta}{2}\right)\norm{\nabla f(x_k)}^2 .
\]

取 \(\eta\le 1/L\)，括号里的因子不小于 \(1/2\)，每一步的函数值确定地下降，除非梯度已经是零。光滑性单独就买到了``往下走''这件事。

凸性原本是在下降引理之后才上场的。它提供的是切平面的全局下界性质：对凸函数，\(f(y)\ge f(x)+\ip{\nabla f(x)}{y-x}\) 对一切 \(y\) 成立，见\hyperref[sec:09-Convex-Optimization/03-Convex-Functions/02]{``切平面为何是全局下界''}。在驻点处梯度为零，这个不等式退化成 \(f(y)\ge f(x)\)——对一切 \(y\)。梯度为零因此直接等价于全局最优。去掉凸性，切平面照样存在，只是它不再托住整条曲线。

\begin{center}
\begin{tikzpicture}[x=1.3cm,y=1.0cm,font=\small,>=stealth]
  \node[font=\footnotesize,black!70] at (0.1,2.55) {凸：切线托住整条曲线};
  \draw[black!75,thick] plot[domain=-1.7:1.7,samples=70] (\x,{0.55*\x*\x-0.3});
  \draw[black!45,dashed] (-1.75,-0.3) -- (1.85,-0.3);
  \fill[black!80] (0,-0.3) circle (1.6pt);
  \node[font=\footnotesize,black!70] at (0.62,-0.62) {\(\nabla f=0\)};
  \node[font=\footnotesize,black!55] at (-1.05,-0.68) {切线在图像下方};
  \begin{scope}[shift={(5.0,0)}]
    \node[font=\footnotesize,black!70] at (0.1,2.55) {非凸：同一条切线漏掉了深谷};
    \draw[black!75,thick] plot[domain=-1.72:1.92,samples=110] (\x,{0.4*\x*\x*\x*\x-1.2*\x*\x-0.45*\x+1.0});
    \draw[black!45,dashed] (-1.8,0.628) -- (2.0,0.628);
    \fill[black!80] (-1.117,0.628) circle (1.6pt);
    \fill[black!80] (1.31,-0.471) circle (1.6pt);
    \draw[black!50] (-1.10,0.55) -- (-0.98,0.20);
    \node[font=\footnotesize,black!70] at (-0.62,0.10) {算法停在这里};
    \draw[<->,black!55] (1.31,0.628) -- (1.31,-0.471);
    \node[font=\footnotesize,black!60] at (1.85,0.1) {漏掉的};
    \node[font=\footnotesize,black!55] at (0.05,-0.75) {曲线有一段跑到切线下方};
  \end{scope}
\end{tikzpicture}

\emph{梯度为零意味着切线是水平的，两边都成立；切线是否为整条曲线的下界，只有左边成立。}
\end{center}

图右边那条虚线和左边那条是同一件事的产物：\(\nabla f=0\)。差别在虚线之下有没有东西。左边没有，所以那个点是全局最优；右边有一整段，所以那个点只是一处平地。算法在停下来的那一刻拿到的信息完全相同——函数值、梯度、也许还有曲率——它无法从中读出自己落在哪一边。

\subsection{累加得到的是最小梯度范数，不是收敛}

即使不知道终点在哪，下降的累加仍能榨出一个定量结论。设 \(f\) 有有限下界 \(f_{\inf}>-\infty\)，取 \(\eta=1/L\)，把每步的下降不等式从 \(k=0\) 加到 \(T-1\)，左侧的相邻项互相抵消，得到

\[
\begin{aligned}
\frac{\eta}{2}\sum_{k=0}^{T-1}\norm{\nabla f(x_k)}^2
&\le f(x_0)-f(x_T)\le f(x_0)-f_{\inf},\\
\min_{0\le k<T}\norm{\nabla f(x_k)}^2
&\le \frac{1}{T}\sum_{k=0}^{T-1}\norm{\nabla f(x_k)}^2
\le \frac{2\bigl(f(x_0)-f_{\inf}\bigr)}{\eta T}.
\end{aligned}
\]

条件对一遍账。第一行用到光滑性（提供每步下降）和下界的存在（提供总落差的天花板），没有用到凸性、没有用到最小值可达、也没有用到可行域紧致。第二行只是``最小值不超过平均值''，纯算术。结论是：前 \(T\) 步里出现过的最小梯度范数按 \(O(1/\sqrt{T})\) 衰减，要把它压到 \(\varepsilon\) 以下，\(O(1/\varepsilon^2)\) 次迭代够用。这是光滑非凸优化的基准复杂度。

请注意主语。定理说的是``存在某个 \(k<T\) 使得那一步的梯度小''，不是``序列收敛''，也不是``最后一步的梯度小''。算法完全可以路过一个梯度极小的点，然后走开，梯度重新变大。实践里之所以感觉不到这个区别，是因为固定步长的梯度下降在光滑函数上通常单调下降，梯度不会大幅反弹；但这是经验而非定理内容。

总落差 \(f(x_0)-f_{\inf}\) 出现在分子上，这一项不能靠调算法改小。初始化得好，它就小，需要的迭代数就少；代价是你从一开始就待在一个不错的区域，本来也没多少可降的。它衡量的是问题给你留了多少下降空间，不是算法的效率。

\subsection{梯度范数小对应三种不同的处境}

上面那把尺子很窄。它认证``脚下平坦''，此外一概不管。这句话在实际训练里会分岔成三种处境，而它们需要完全不同的对策。

第一种是真的到了局部极小。Hessian 正定，附近确实没有更低的点，继续训练只是在消耗预算。

第二种是停在鞍点或近似鞍点上。梯度为零，一阶条件完全通过，但存在负曲率方向，沿它走函数值还能降。高维参数空间里这种处境远比第一种常见，怎么把它和第一种区分开、又怎么逃离，是\hyperref[sec:09-Convex-Optimization/08-Nonconvex-and-Stochastic-Optimization/02]{``二阶结构区分局部极小与鞍点''}的内容。

第三种最容易被误读：梯度在变小，但不是因为接近了什么解，而是因为参数正沿着一条无限延伸的缓坡跑向无穷。一维的 \(f(x)=\arctan x\) 就够说明问题。它光滑，全实线上可取 \(L=1\)，下确界 \(-\pi/2\) 有限，导数 \(f'(x)=1/(1+x^2)\) 在 \(x\to-\infty\) 时趋于零。梯度下降跑到 \(x\approx-300\) 时 \(f'(x)\approx1.1\times10^{-5}\)，停止条件 \(\norm{\nabla f}<10^{-5}\) 眼看就要触发，此时函数值离下确界还差 \(0.004\)。差得不算多，但下确界不在任何有限点被取到——你没有一个可以返回、可以部署的参数向量，只有一个越跑越远的方向。

这个现象在实际模型里有精确对应。线性可分数据上不加正则的 logistic 回归，损失的下确界是零，达到零要求所有样本的 margin 趋于正无穷，也就是 \(\norm{w}\to\infty\)。训练日志上会同时出现三条线索：损失单调下降、梯度范数单调下降、参数范数单调上升。前两条单看都像收敛，第三条才暴露真相。识别办法很直接：把 \(\norm{w}\) 也记进日志，它若在长时间里近似按 \(\log t\) 增长，就该加 \(L_2\) 正则或提前停止，而不是继续调步长。

在这三种处境里，梯度范数衡量的都是当前地形的平坦程度，而不是离一个有限解的距离。用它单独认证收敛，是拿本地的量去回答全局的问题。

\subsection{局部光滑性仍然可用，全局常数不可用}

理论里的 \(L\) 在深度网络上没有可用的数值。同一个损失曲面上，不同区域的曲率能差好几个数量级；同一点处，不同层的梯度尺度也能差好几个数量级。按最坏情形取 \(L\) 会得到一个小到无法训练的步长。

但下降引理的推导只在当前点和试探点之间的一小段线段上用到 Lipschitz 控制，那一段上的局部常数就够了。回溯线搜索正是为此设计的：从上一次用过的步长起步，检查 Armijo 条件

\[
f(x_k-\eta\nabla f(x_k))\le f(x_k)-c\,\eta\norm{\nabla f(x_k)}^2 ,\qquad 0<c<1 ,
\]

不满足就把 \(\eta\) 乘上一个收缩因子重试。整个过程只用当前点的函数值、梯度和试探点的函数值，这三样东西在非凸情形下同样拿得到。

要说清的是这里也发生了一次降级。凸问题里线搜索找到的下降量能通过切平面不等式换算成到最优解的距离缩小；非凸问题里它只认证``这一步的函数值确实降了''，认证范围随这一步结束而结束，下一步得重新认证。所以线搜索在非凸训练里的作用是保证数值稳定和避免发散，不是保证进展。

停止条件也该相应地放宽口径。梯度范数只是候选信号之一，配上参数的相对变化量、最近若干步损失的下降幅度、以及验证集风险是否还在改善，一起看才有意义。它们回答的问题并不相同：前两个说的是迭代是否稳定，第三个说的是优化是否还有余地，最后一个说的是继续优化是否还值得——最后这条与训练损失的关系是间接的，留到\hyperref[sec:09-Convex-Optimization/08-Nonconvex-and-Stochastic-Optimization/06]{``优化误差、统计误差与泛化不能混为一谈''}再拆。

眼下更急的问题是第一种和第二种处境怎么分。梯度范数在两者上读数一样，要区分它们，得往上再要一阶导数。
